%%%%%%%%%%%%%%%%%%%%%%%%%%%%%%%%%%%%%%%%
%% Other packages
%%%%%%%%%%%%%%%%%%%%%%%%%%%%%%%%%%%%%%%%
\usepackage{booktabs}
\usepackage[detect-mode=true]{siunitx}
\usepackage{wrapfig}
\usepackage{xcolor}
\usepackage{graphicx}
\graphicspath{{./figures/}}
\usepackage{amsmath}
\usepackage{amsfonts}
\usepackage{cleveref}
\usepackage{vector} % local
\usepackage[font=small]{caption}
\usepackage{natbib}
\usepackage{subcaption}
\usepackage{adjustbox}
\usepackage{longtable}
\usepackage[flushleft]{threeparttable}
\usepackage{multirow}
\usepackage{dsfont} % \mathds{1}
\usepackage{colortbl}
\usepackage[mode=buildnew]{standalone}
% \usepackage[natbib=true]{biblatex}
\usepackage{tabularx}


\newcommand{\algname}{\text{ARIS}}
\newcommand{\algfullname}{\text{Amortized Rare-Event Importance Sampling}}

\sisetup{separate-uncertainty=true,multi-part-units=single,table-align-uncertainty=true, table-format = -3.2(3), detect-weight=true, detect-family=true, detect-all, input-digits = 0123456789\pi}


\definecolor{juliafunccolor}{HTML}{2e51a2}
\definecolor{juliaparenscolor}{HTML}{f2d71c}
\definecolor{pastelRed}{RGB}{245, 97, 92}
\definecolor{pastelBlue}{RGB}{0, 114, 178}
\definecolor{pastelGreen}{RGB}{0, 158, 115}
\definecolor{pastelPurple}{RGB}{135, 112, 254}
\definecolor{pastelYellow}{RGB}{255, 205, 117}


\definecolor{newcolor}{HTML}{000000}
\newcommand{\new}[1]{{\color{newcolor}#1}}

\newcommand\blfootnote[1]{%
  \begingroup
  \renewcommand\thefootnote{}\footnote{#1}%
  \addtocounter{footnote}{-1}%
  \endgroup
}

%%%%%%%%%%%%%%%%%%%%%%%%%%%%%%%%%%%%%%%%
% Algorithms
%%%%%%%%%%%%%%%%%%%%%%%%%%%%%%%%%%%%%%%%
\usepackage{algorithmicx}
\usepackage{algorithm}
\usepackage{algpseudocode}


%%%%%%%%%%%%%%%%%%%%%%%%%%%%%%%%%%%%%%%%
%% TiKZ and PGF plots
%%%%%%%%%%%%%%%%%%%%%%%%%%%%%%%%%%%%%%%%
\usepackage{tikz}
\usepackage{pgfplots}
\pgfplotsset{compat=newest}


\usetikzlibrary{calc}
\usetikzlibrary{shapes.geometric}
\usetikzlibrary{external}
\usetikzlibrary{patterns}
\usetikzlibrary{shapes,arrows,fit}
\usetikzlibrary{positioning}
\usetikzlibrary{arrows.meta, calc, shapes}
\usetikzlibrary{graphs}
\usetikzlibrary{decorations.pathmorphing}
\usetikzlibrary{decorations.pathreplacing}
\usetikzlibrary{backgrounds}
\usetikzlibrary{shadows}
\usepgfplotslibrary{fillbetween}

\definecolor{blues1}{RGB}{198, 219, 239}
\definecolor{blues2}{RGB}{158, 202, 225}
\definecolor{blues3}{RGB}{107, 174, 214}
\definecolor{blues4}{RGB}{49, 130, 189}
\definecolor{blues5}{RGB}{8, 81, 156}

\definecolor{grays1}{RGB}{219, 219, 219}
\definecolor{grays2}{RGB}{202, 202, 202}
\definecolor{grays3}{RGB}{174, 174, 174}
\definecolor{grays4}{RGB}{130, 130, 130}
\definecolor{grays5}{RGB}{81, 81, 81}

%% For Pendulum
\usetikzlibrary{calc,tikzmark,positioning,shadows,arrows}

\definecolor{pastelred}{RGB}{245,97,92}
\definecolor{steelblue}{rgb}{0.27,0.51,0.71}

% Format Helpers
\newcommand{\ph}[1]{{\textbf{#1}:}} %
\newcommand{\red}[1]{\textcolor{red}{#1}}

%% Math Helpers
\newcommand{\ones}{\mathbf 1}
\newcommand{\reals}{{\mbox{\bf R}}}
\newcommand{\integers}{{\mbox{\bf Z}}}
\newcommand{\symm}{{\mbox{\bf S}}}  % symmetric matrices

\newcommand{\nullspace}{{\mathcal N}}
\newcommand{\range}{{\mathcal R}}
\newcommand{\Rank}{\mathop{\bf Rank}}
\newcommand{\Tr}{\mathop{\bf Tr}}
\newcommand{\diag}{\mathop{\bf diag}}
\newcommand{\card}{\mathop{\bf card}}
\newcommand{\rank}{\mathop{\bf rank}}
\newcommand{\conv}{\mathop{\bf conv}}
\newcommand{\prox}{\mathbf{prox}}

\newcommand{\Expect}{\mathop{\mathbb{E}}}
\newcommand{\Reals}{\mathop{\mathbb{R}}}
\newcommand{\Prob}{\mathop{\bf Prob}}
\newcommand{\Co}{{\mathop {\bf Co}}} % convex hull
\newcommand{\dist}{\mathop{\bf dist{}}}
\DeclareMathOperator*{\argmin}{\arg\!\min}
\newcommand{\epi}{\mathop{\bf epi}} % epigraph
\newcommand{\Vol}{\mathop{\bf vol}}
\newcommand{\dom}{\mathop{\bf dom}} % domain
\newcommand{\intr}{\mathop{\bf int}}
\newcommand{\sign}{\mathop{\bf sign}}

\newcommand{\cf}{{\it cf.}}
\newcommand{\eg}{{\it e.g.}}
\newcommand{\ie}{{\it i.e.}}
\newcommand{\etc}{{\it etc.}}

\newcommand{\mbf}[1]{\boldsymbol{#1}}

\newcommand{\bx}{\mathbf{x}}
\newcommand{\by}{\mathbf{y}}
\newcommand{\bz}{\mathbf{z}}

\usetikzlibrary{pgfplots.colorbrewer}
\usepgfplotslibrary{groupplots}
\usepgfplotslibrary{polar}
\usepgfplotslibrary{smithchart}
\usepgfplotslibrary{statistics}
\usepgfplotslibrary{dateplot}
\usepgfplotslibrary{ternary}

\pgfplotsset{
    colormap={pastelPG}{
        color(0cm)=(pastelPurple)   % Start color: blue
        color(0.5cm)=(white) % Neutral color: white
        color(1cm)=(pastelGreen)    % End color: red
    }
}